\subsection{BioseqDB all cache layers}
\label{sec:pg-storage-caches}

The optional column cache stores decompressed stripe chunks. It is separate
from the scanner-owned packed fragment slots in
Figure~\ref{fig:sequence-scan-lifecycle}.

The setting \path{bioseq.column_cache.enabled} controls the decompressed cache.
The original benchmark campaign disabled it, whereas binary releases enable
it by default. It stores decoded chunks inside one PostgreSQL backend process.
Its key contains relation, stripe, chunk
group, and column identifiers (the configured
capacity when enabled defaults to 200~MiB)

The cache can return borrowed pointer to a decompression buffer for a chunk, which allows to prevent additional chunk loads wasting time on decompression multiple times.
As I discussed all methods of caching within BisoeqDB already, I present table with comparison of all cache layers:

\begin{table}[htbp]
\centering
\small
\caption{Different types of cahces relevant in context of BioseqDB}
\label{tab:storage-cache-scopes}
\begin{tabularx}{\linewidth}{@{}p{.23\linewidth}XX@{}}
\toprule
\textbf{Name of the cache method} & \textbf{Saved data} & \textbf{Description} \\
\midrule
PostgreSQL buffer & Relation pages & Internal buffer manager in Postgres loads/unloads pages and can pin pages to optimise page access. In BioseqDB I don't go to such low level of abstraction and assume page loading is opaque mechanism. \\
Column cache~\ref{sec:pg-storage-caches} & Buffer with decompressed column chunks & Can cache decompressed chunks, so loading and decompression of the same chunks can be efficient. \\
Fragment slots~\ref{fig:sequence-scan-lifecycle} & Copy of one packed fragment & Each scanner can store fragments in its own memory, so when pipeline needs to lookup sequence parts not far away it will resuse already loaded fragments \\
Query-result cache~\ref{sec:seq-results-cache} & Completed SQL result rows & Separate \path{bioseq_cache} subextension can store and share results after pipeline execution is finished. This extenion is disabled by default, but it can read saved results if the queries are same. \\
\bottomrule
\end{tabularx}
\end{table}
